\section{边疆与南方：王室的地图走到哪里，能力就走到哪里}

西周的政治重心在关中和黄河中下游，但王室并不满足于核心区。向南的军事行动、与周边政治共同体的关系以及交通补给的难度，提醒我们：王权不是一条画在地图上的边界，而是补给、沟通、盟友和军队能否持续抵达的范围。

\begin{event}{昭王在位晚期（约前10世纪末）}{周昭王南征失利}{王年推定}{汉水流域}\eventanchor{WZ-0977-ZHAONAN}{西周·昭王南征}
\term{这一年发生了什么}：传世材料记载周昭王南征，“丧六师于汉”，王本人也未能返回。昭王在位常被推定为约前977至前957年，南征具体年份没有金文直接标明；但王室在南方遭遇重大军事损失这一点具有稳定的叙事核心。

\term{事情为什么走到这里}：南方资源、交通和政治关系吸引周王室继续扩张；但汉水流域并不是关中周人可以轻易套用旧式动员的地方。远征意味着漫长补给线、陌生水土和难以预测的地方联盟。

\term{当事人的选择与代价}：向外用兵可能带来资源与声望，也可能暴露王室的后勤极限。昭王之死使这种代价不再抽象：王位、军队和威望同时受损。

\term{后来改变了什么}：\eventref{WZ-0977-ZHAONAN}{西周·昭王南征}不能单独解释西周衰落，却让我们看到王室资源并非无限。后续王室若要同时维持东方网络、西部边防和内部赏赐，财政与动员压力会越来越大。
\end{event}

\begin{causalchain}[边疆不是背景板]
西周向南的挫败说明，王室的军事能力必须和地理、补给、地方政治一起计算。边疆的高成本会回流到中心：军队损失要补，赏赐要兑现，继承与贵族关系也更难维持。
\end{causalchain}

\begin{sourcenote}[这场失败能知道多少]
\sourcebook{史记·周本纪}只简记“昭王南征不复”，\sourcebook{竹书纪年}称“丧六师于汉”，《左传》还有楚人对周使的讽刺性追述。小盂鼎和湖北地区西周铜器可以提示南方军事与文化渗透，但不能直接证明昭王在何地、何年败亡。铜锡资源、宣示王权和南方政治关系都是合理解释，不能选其中一项当作已经证实的唯一动机。
\end{sourcenote}
